\documentclass[12pt]{article}
\usepackage[utf8]{inputenc}
\usepackage[T1]{fontenc}
\usepackage{hyperref}
\title{Unami Language Resources for K-12 Curriculum Development}
\author{Compiled by Claude Sonnet 5 (claude-sonnet-5) for Andr\'as Kornai}
\date{}
\begin{document}
\maketitle
\section{Introduction}
Unami (ISO 639-3 unm) is a dormant Delaware (Lenape) language, now used mainly as a second language. This report lists only resources that appeared in web search results; the search surfaced little beyond the Delaware Tribe's own materials.
\section{Primary Educational Resources}
The Lenape Talking Dictionary has a lessons section, but the lesson pages themselves were not visible in the results, so their content is unverified. \url{https://talk-lenape.org/}
\section{Linguistic and Reference Materials}
The Lenape Talking Dictionary is described as the official dictionary of the Delaware Tribe's language. It is searchable in English and Lenape, and entries carry speaker audio. Reported size is inconsistent (about 15,000 to over 17,000 words). \url{https://talk-lenape.org/} A sample entry: \url{https://talk-lenape.org/detail?id=15664}
\section{Religious and Community Resources}
The Delaware Tribe describes its Lenape Language Preservation Project here: \url{https://delawaretribe.org/cultural-education/culture-and-language/lenape-language/} and here: \url{https://delawaretribe.org/?p=6159}
\section{Digital and Audio Resources}
The dictionary includes audio, some of it from recordings made as early as 1928. The directory native-languages.org lists audio phrasebooks drawn from Nora Thompson Dean's Lenape cassettes. \url{https://www.native-languages.org/lenape.htm}
\section{Curriculum Development Context}
No school curriculum, textbook or primer was found in these searches. The dictionary and the tribal project are the obvious starting points. Contact the Delaware Tribe directly about any teaching materials.
\end{document}